% !TeX root = ../机器学习笔记.tex
% 由主文档通过 \input 载入；请编译主文档。
\setcounter{chapter}{17} % 保留原书章号。
\chapter{深度模型与表示学习初步}
\label{chap:ml-representation}

表示学习（representation learning）把用于预测的特征也纳入训练。卷积、循环和注意力并非互不相关的技巧，而是对数据局部性、顺序和交互方式的不同假设。

\section{表示、归纳偏置与特征学习}

表示 $z=\phi_\theta(x)$ 经预测头 $g_\psi(z)$ 得到输出，两部分可联合训练。若变换 $T$ 后表示不变，即 $\phi(Tx)=\phi(x)$，称具有不变性（invariance）；若存在输出变换 $T'$ 使 $\phi(Tx)=T'\phi(x)$，称具有等变性（equivariance）。识别图像类别可能希望对小平移近似不变，定位物体则应保留位置变化。

表示是否有用依赖任务。把全部输入映射到同一个常数能实现很多不变性，却丢失预测信息；重构输入很好也可能保留大量与标签无关的细节。评价表示时应明确下游标签量、预测头能力和是否允许微调。

\section{卷积网络与局部结构}

卷积神经网络（convolutional neural network，CNN）使用局部连接和跨位置共享权重。常见一维运算为 $z_t=\sum_{j=0}^{k-1}w_jx_{ts+j-p}+b$，其中 $k$ 为核宽、$s$ 为步幅、$p$ 为两端填充长度；严格说这是互相关，深度学习中常称卷积。长度 $T$ 的输入在合法参数下产生 $\lfloor(T+2p-k)/s\rfloor+1$ 个输出。

共享权重把各位置的同类局部模式联系起来，池化进一步压缩位置分辨率。多层堆叠扩大感受野，但有限边界、填充与下采样会影响严格平移等变性。卷积表达局部结构先验，并不保证其适合任意表格数据。

\section{循环网络与序列依赖}

循环神经网络（recurrent neural network，RNN）用 $h_t=\phi(W_hh_{t-1}+W_xx_t+b)$ 压缩过去输入，输出可在每步产生或仅由最终状态产生。随时间反向传播就是展开递推后的链式法则；长序列中的 Jacobian 连乘解释了梯度消失与爆炸，门控机制提供较易控制的信息保留路径。

序列预测必须规定每个时刻可用的信息。双向网络可利用完整输入序列，却不适用于要求只看过去的在线预测。教师强制（teacher forcing）在自回归训练中提供真实前序输出，推断时可能需要使用模型自身输出，两者产生的历史分布可能不同；截断反向传播则改变梯度传播的时间范围。

\section{注意力机制与 Transformer}

对查询矩阵 $Q\in\mathbb R^{T_q\times d_k}$、键 $K\in\mathbb R^{T_k\times d_k}$、值 $V\in\mathbb R^{T_k\times d_v}$，缩放点积注意力为
\begin{equation}
 \operatorname{Att}(Q,K,V)=
 \operatorname{softmax}(QK^\top/\sqrt{d_k}+M)V.
 \label{eq:ml-attention}
\end{equation}
Softmax 逐行计算，允许位置的掩码 $M$ 为零，不允许位置视为负无穷，每行至少保留一个位置。每个输出因此是值向量的凸组合，权重由查询和键的相容性决定。

自注意力从同一序列生成三组向量；多头注意力学习多组投影后合并。Transformer 还使用逐位置前馈层、残差与归一化，并用位置表示或其他机制提供顺序信息。因果掩码禁止访问未来位置。长度为 $T$ 的标准全注意力计算成对分数需 $O(T^2d_k)$ 运算，显式存储分数需 $O(T^2)$ 空间；架构效率必须结合实际算法与维度评价。

\section{自编码、自监督与迁移表示}

自编码器（autoencoder）由编码器 $z=e_\theta(x)$ 与解码器 $\widehat x=d_\psi(z)$ 组成，通过重构损失学习表示。没有容量约束时，记忆或恒等映射可能使重构很好却缺乏有用压缩。去噪、瓶颈或稀疏约束分别加入不同偏好。

自监督任务可预测被遮蔽的内容，或让同一样本的两种增强表示相近。以相似度 $s$ 和温度 $\tau>0$ 为例，对比损失可写为
$-\log[e^{s(z,z^+)/\tau}/\sum_{u\in\mathcal C}e^{s(z,u)/\tau}]$，
候选集 $\mathcal C$ 包含正例和所选负例。增强若改变标签语义，或负例包含语义上相同对象，训练目标可能偏离下游任务。自回归生成以条件概率链式分解训练，变分自编码器则引入隐变量并优化 ELBO；二者的生成目标也不自动保证表示对所有任务有用。

\section{例题、练习与学习检查}

\begin{example}[注意力的加权平均]
一个查询对应两个键，缩放分数为 $0,\log3$，值为标量 $2,6$，均允许访问。
\end{example}
\begin{solve}
Softmax 权重为 $1/4,3/4$，输出为五。若第二个位置是未来位置并被因果掩码屏蔽，权重变为 $1,0$，输出为二。掩码改变的是信息可用范围，并非简单调整一个评价阈值。
\end{solve}

习题：用核 $(1,-1)$ 对输入 $(1,2,3,4)$ 做无填充、单位步幅互相关；检验一个序列模型是否使用未来信息；构造会改变标签的图像或文本增强；比较固定随机表示与训练表示时，控制下游模型和标签预算。




